\section{Параллельная обработка}
\label{sec:parallelism}

Pull выбран в разделе~\ref{sec:approaches}. Здесь --- как этот проход распараллелить, не получив гонок
и не потеряв детерминизм.

\subsection{Без блокировок и детерминизм}
Каждый поток берёт непрерывный диапазон приёмников и пишет только в свой непересекающийся кусок
$s_{\mathrm{new}}$. Массивы $\mathrm{contrib}$ и $s_{\mathrm{old}}$ общие и только на чтение. Отсюда
два следствия.
\begin{itemize}[itemsep=3pt]
  \item \textbf{Горячий цикл без блокировок.} В проходе $\bigO(m)$ нет ни атомиков, ни
        \texttt{synchronized}. Единственная синхронизация --- один барьер на итерацию.
  \item \textbf{Детерминизм, не зависящий от числа потоков.} Никакие два потока не пишут в одну
        ячейку, а порядок сложения внутри ячейки задан порядком рёбер в файле. Поэтому результат
        получается бит в бит одинаковым при $1$ и при $24$ потоках и от запуска к запуску. Планировщик
        потоков на ответ не влияет: ни одна ячейка не делится между потоками, а порядок суммирования в
        ней зафиксирован на диске.
\end{itemize}

\subsection{Балансировка и гипер-узлы}
Делить поровну по числу вершин нельзя. Из-за степенного распределения равные по числу вершин куски
содержат очень разное число рёбер, и поток с гипер-узлом становится отстающим. Поэтому:
\begin{itemize}[nosep]
  \item границы кусков выбираются так, чтобы уравнять число рёбер (по префиксной сумме входящих
        степеней).
  \item кусков создаётся больше, чем потоков (в несколько раз), и пул разбирает их по мере
        освобождения. Это сглаживает остаточный перекос.
  \item гипер-узел в модели pull неделим (у ячейки один владелец), но в абсолютных числах он мал.
        Входящая степень $50\,000$ --- это около $0{,}2$~МБ источников, микросекунды работы.
\end{itemize}

\subsection{Реализация на Java}
\label{sec:java}
\begin{itemize}[itemsep=3pt]
  \item \textbf{Пул потоков и барьер.} Фиксированный \texttt{ExecutorService} на $P$ потоков по числу
        ядер. На каждой итерации задачи-куски отдаются в \texttt{invokeAll}, который ждёт завершения
        всех. Это ожидание и есть барьер итерации: пока все приёмники не пересчитаны, следующая
        итерация не начнётся. Отдельный \texttt{CyclicBarrier} не нужен.
  \item \textbf{Корректность по модели памяти Java.} Потоки пишут в $s_{\mathrm{new}}$ без
        \texttt{volatile} и без блокировок. Это безопасно, потому что завершение задачи и его
        наблюдение через \texttt{invokeAll} задают отношение happens-before. Всё, что поток записал,
        видно главному потоку, когда тот увидел задачу завершённой. Значит записи итерации $t$
        опубликованы до чтений итерации $t{+}1$.
  \item \textbf{Разбиение диапазонами против false sharing.} Поток пишет непрерывный кусок, поэтому
        делиться с соседом может только одна кэш-линия на границе. При кусках в тысячи элементов это
        незаметно. Чередующееся разбиение, наоборот, посадило бы соседние потоки на общие кэш-линии.
  \item \textbf{Детерминированные суммы.} Глобальные суммы (масса фонового узла и норма сходимости)
        считаются в фиксированном порядке: каждый поток отдаёт частичную сумму, а главный поток
        складывает их по порядку номеров. \texttt{DoubleAdder} и \texttt{stream().sum()} с
        неопределённым порядком сознательно не используются. Потоки платформенные, а не виртуальные:
        виртуальные хороши для блокирующего ввода-вывода, но не ускоряют счёт.
  \item \textbf{Мало мусора.} Массивы $s_{\mathrm{old}}$, $s_{\mathrm{new}}$, $\mathrm{contrib}$,
        \texttt{outdeg}, \texttt{sourcesPtr} выделяются один раз. Между итерациями $s_{\mathrm{old}}$
        и $s_{\mathrm{new}}$ просто меняются ссылками. Горячий цикл почти не создаёт мусора, поэтому
        Serial GC справляется даже при крошечной куче.
\end{itemize}

\subsection{Ввод-вывод}
Граф не помещается в память, поэтому файл \texttt{sources} перечитывается с диска на каждой итерации.
Закэшировать его целиком нельзя, в этом суть задачи. Каждый поток читает свои куски последовательно,
держа только буфер фиксированного размера ($64$~КиБ), а не весь кусок. Буфер свой на каждый поток и
переиспользуется между итерациями, поэтому суммарная память на буферы --- это $P\cdot B$, а не $P$
копий чего-то большого. На SSD упреждающее чтение ОС неплохо совмещает ввод-вывод со счётом.
